\sectionTitle{Ausgewählte Forschungsarbeiten}{\faTasks}

\begin{projects}
    \project
	{Rana, A., Hung, C.C., Sun, Q., Kunkel, J.M. and Lawrence, C., 2026. Oblivion: Self-Adaptive Agentic Memory Control through Decay-Driven Activation. arXiv preprint arXiv:2604.00131.}{\textit{Preprint auf arXiv verfügbar, ACL ARR Einreichung Mai 2026}}
	{
	\website{https://arxiv.org/abs/2604.00131}{arxiv.org/abs/2604.00131}
	\website{https://nec-research.github.io/oblivion/}{nec-research.github.io/oblivion/}
	}
	{Entwicklung eines kognitiv inspirierten Frameworks zur Gedächtnissteuerung, das Vergessen als zerfallsbedingte Verringerung der Zugänglichkeit und nicht als explizites Löschen auffasst.}
	{Agentisches Gedächtnis, Langzeitgedächtnis, Gedächtnisverwaltung}

    \project
	{Rana, A., Shaker, A., Saralajew, S., Suzuki, T., Yasuda, K., Kato, S., Wada, T., Fujikawa, T. and Kikutsuji, T. (2025) 'Prototype-Based Learning for Healthcare: A Demonstration of Interpretable AI', in Proceedings of the Demo Track of the IEEE International Conference on Data Mining Workshops 2025 (ICDMW 2025).}{\textit{Best Demo Paper Award, ICDM Demo Track 2025}}
	{
	\website{https://ieeexplore.ieee.org/document/11415757}{10.1109/ICDMW69685.2025.00323}
	\website{https://arxiv.org/abs/2601.02106}{arxiv.org/abs/2601.02106}
	}
	{Entwicklung eines Systems auf dem Stand der Technik für interpretierbare Krankheitsrisikoprognosen und Erklärungen; zusätzlich Einsatz prototypenspezifischer Autoencoder zur Modellierung von Abhängigkeiten zwischen Merkmalen für Interventionskorrekturen.}
	{Prototypenbasiertes Lernen, Interventionen und Empfehlungen}

	\project
	{Saralajew, S., Rana, A., Villmann, T. and Shaker, A. (2025) 'A Robust Prototype-Based Network with Interpretable RBF Classifier Foundations', in Proceedings of the AAAI Conference on Artificial Intelligence, Vol. 39, No. 19, pp. 20273-20282.}{\textit{AAAI 2025}}
	{
	\website{https://github.com/nec-research/enlaight}{github.com/nec-research/enlaight}
	\website{https://github.com/si-cim/cbc-aaai-2025}{github.com/si-cim/cbc-aaai-2025}
	\website{https://arxiv.org/abs/2412.15499}{arxiv.org/abs/2412.15499}
	}
	{Entwicklung tiefer Classification-by-Components-Architekturen, Durchführung umfassender Interpretierbarkeitsanalysen und Entwicklung von Software-Toolkits zur Reproduzierbarkeit.}
	{Classification-by-Components, Prototypenbasiertes Lernen, Interpretierbarkeits-Benchmark auf dem Stand der Technik}

	\project
	{Rana, A., Oesterle, M. and Brinkmann, J. (2024) 'GOV-REK: Governed Reward Engineering Kernels for Designing Robust Multi-Agent Reinforcement Learning Systems', in Proceedings of the 23rd International Conference on Autonomous Agents and Multiagent Systems (AAMAS '24). Richland, SC: International Foundation for Autonomous Agents and Multiagent Systems, pp. 2429-2431.}{\textit{Extended Abstract, AAMAS 2024}}
	{
	\website{https://github.com/arana-initiatives/gov-rek-marls}{github.com/arana-initiatives/gov-rek-marls}
	\website{https://arxiv.org/abs/2404.01131}{arxiv.org/abs/2404.01131}
	}
	{Verwendung von Priors für Belohnungsverteilungen und des Hyperband-Algorithmus, um Belohnungen für Agenten in dünn belohnten Multi-Agenten-Systemen (MAS) effizient und automatisiert zu entwickeln und schnelleres Lernen zu ermöglichen.}
	{Zentralisiertes und dezentrales MAS-Training, Gaußsche Kernelverteilungen, Training und Hyperparameteroptimierung}

	\project
	{Rana, A., Golchha, P., Juntunen, R., Coajă, A., Elzamarany, A., Hung, C.C. and Ponzetto, S.P. (2022) 'LeviRANK: Limited query expansion with voting integration for document retrieval and ranking: notebook for the touché lab on argument retrieval at CLEF 2022', in CEUR Workshop Proceedings, Vol. 3180, pp. 3074-3089. RWTH Aachen.}{\textit{Touch\'e IR Workshop, CLEF 2022}}
	{
	\website{https://github.com/softgitron/LeviRank}{github.com/softgitron/LeviRank}
	\website{https://ceur-ws.org/Vol-3180/paper-259.pdf}{https://ceur-ws.org/Vol-3180/paper-259.pdf}
	}
	{Entwicklung eines neuartigen, zur Abfragerelevanz proportionalen Retrieval-Ansatzes und Implementierung des Designmusters ``Expando-Mono-Duo'' mit T5-Modellen.}
	{BM25- und semantisches Retrieval, Relevanzfeedback, Punktweises und paarweises Ranking}

	\project
	{Rana, A., Khanna, D., Ghosal, T., Singh, M., Singh, H. and Rana, P.S. (2022) 'RerrFact: Reduced Evidence Retrieval Representations for Scientific Claim Verification', in Proceedings of the Scientific Document Understanding Workshop at AAAI 2022. CEUR Workshop Proceedings, Vol. 3164.}{\textit{SDU Workshop, AAAI 2022}}
	{
	\website{https://github.com/ashishrana160796/rerrfact}{github.com/RerrFact}
	\website{https://arxiv.org/abs/2202.02646}{arxiv.org/abs/2202.02646}
	}
	{Neugestaltung mehrstufiger Information-Retrieval- und Klassifikationsmodule mit lose gekoppeltem Modelltraining über verschiedene Stufen hinweg, um maximalen Präzisionsdurchsatz zu erreichen.}
	{Überprüfung wissenschaftlicher Behauptungen, Modulares Pipeline-Design, Transformer-Sprachmodelle}

	\project
	{Rana, A., Singh, H., Mavuduru, R., Pattanaik, S. and Rana, P.S., 2022. Quantifying prognosis severity of COVID-19 patients from deep learning based analysis of CT chest images. Multimedia Tools and Applications, 81(13), pp.18129-18153.}{\textit{MTAP Journal 2022}}
	{
	\website{https://doi.org/10.1007/s11042-022-12214-6}{doi.org/10.1007/s11042-022-12214-6}
	}
	{Entwicklung einer Single-Shot-Detector-(SSD)-Pipeline zur besseren Prognose des COVID-19-Schweregrads anhand eines Datensatzes von CT-Schichtbildern. \textit{(DOI: 10.1007/S11042-022-12214-6)}}
	{Single Shot Detector (SSD), Bildähnlichkeitsnetze, Mehrklassenklassifikation}

	\project
	{Rana, A. and Malhi, A. (2021) 'Building safer autonomous agents by leveraging risky driving behavior knowledge', in 2021 International Conference on Communications, Computing, Cybersecurity, and Informatics (CCCI). IEEE, pp. 1-6.}{\textit{Informatics Best Paper Award, CCCI 2021}}
	{
	\website{https://ieeexplore.ieee.org/document/9583209}{https://ieeexplore.ieee.org/document/9583209}
	}
	{Verbesserung und kausale Bewertung der Agentenleistung durch Nutzung von Informationen über Kollisionsszenarien in einem kausalen Versuchsdesign \textit{(DOI: 10.1109/CCCI52664.2021.9583209)}.}
	{Modellfreies Deep Reinforcement Learning, Stochastische Regelung, Kausales Versuchsdesign}

	\iffalse
	\project
	{Automating Couchbase Monitoring Integrations with Prometheus \& Grafana}{\textit{Couchbase Connect.ONLINE 2020}}
	{\website{github.com/ashishrana160796/cautious-integration-analysis}{GitHub-Repository}}
	{Präsentation eines orchestrierten Multi-Tool-Setups mit Prometheus, Couchbase-Exporter und Grafana für das Monitoring von Couchbase-Datenbanken mit REST-HTTP-Servern sowie hochverfügbarer Architekturen und prädiktiver Warnungen.}
	{REST, Couchbase, Prognosen und Ausreißererkennung}
	\project
	{Exploring Numerical Calculations with CalcNet}{ICTAI 2019}
	{\website{https://ieeexplore.ieee.org/abstract/document/8995315}{IEEE-Portal} \website{github.com/ashishrana160796/calcnet-experiments}{GitHub-Repository}}
	{Erstautor; Vorschlag einer neuronalen Architektur, die mathematische Ausdrücke mit NALU-Schichten parst und approximiert.}
	{NALU, MLP, Numerische Analyse}
	\project
	{Systematically designing better instance counting models on cell images with NALU}{\textit{arXiv.org}}
	{\website{https://arxiv.org/abs/2004.06674}{arXiv: 2004.06674} \website{https://github.com/ashishrana160796/nalu-cell-counting}{GitHub-Repository}}
	{Erstautor; Verbesserung der Genauigkeit auf Zellzähldatensätzen durch NALU-Rückkopplungen in FCRN und U-Nets.}
	{FCRN, U-Net, NALU}
	\fi
\end{projects}
